\part{Implementación}
\label{part:implementacion}
%%%%%%%%%%%%%%%%%%%%%%%%%%%%%%%%%%%%%%%%%%%%%%%%%%%%%%%%%%%%%%%%%

\section{Implementación del sistema}
\label{sec:implementacion-sistema}

\subsection{Estructura del repositorio y componentes}

\textit{Sección por completar a medida que el código estabilice su estructura definitiva. Placeholder estructural con lo que ya es firme.}

El código del TFG se organiza en torno a cuatro directorios principales: \texttt{distill/} (scripts de destilación), \texttt{serving/} (configuraciones y launchers de vLLM), \texttt{eval/} (cliente de evaluación y scripts de post-proceso) y \texttt{routing/} (implementación de las Políticas B y C y el predictor). Los experimentos reproducibles viven en \texttt{experiments/}, cada uno con un fichero de configuración, un script SLURM, y un directorio de resultados trazable al JobID que lo generó.

Los detalles técnicos de cada componente (librerías concretas, versiones, estructura de las configuraciones, formato de logs) se desarrollarán en este capítulo cuando el código del TFG esté estabilizado, para no tener que reescribirlo dos veces si aún hay decisiones abiertas.

\subsection{Integración con SLURM y entorno de ejecución}

\textit{Por completar con scripts SLURM finales.}

Esta sección describirá cómo se despliega el entorno en MareNostrum~5, incluyendo la estrategia de módulos, el entorno Python, los scripts SLURM tipo para destilación, serving y evaluación, y los mecanismos de registro automático de consumo energético vía \texttt{sacct}.

\subsection{Problemas técnicos relevantes y soluciones}

\textit{Se irá completando con los hallazgos reales a medida que aparezcan.}

Este bloque recoge los problemas no triviales encontrados durante la implementación (configuraciones de vLLM, versiones incompatibles, límites de la cola, timeouts, etc.) y cómo se resolvieron, con el objetivo de servir de guía para quien continúe esta línea.

%%%%%%%%%%%%%%%%%%%%%%%%%%%%%%%%%%%%%%%%%%%%%%%%%%%%%%%%%%%%%%%%%
